\documentclass[10pt,a4paper]{amsart}
\usepackage[margin=19mm]{geometry}
\usepackage{amsmath,amssymb,mathtools}
\usepackage[scr=rsfs,cal=euler]{mathalfa}
\usepackage{xfrac}
\usepackage[unicode,hidelinks,bookmarks=false]{hyperref}
\newcommand{\ie}{i.e.\ }
\newcommand{\cf}{cf.\ }
\newcommand{\resp}{respectively\ }
\newcommand{\Subsection}{\subsection}
\providecommand{\nobreakdash}{\nobreak\hbox{-}}
\setlength{\emergencystretch}{2em}
\multlinegap0pt
\usepackage{bm}
\usepackage{bbm}
\usepackage{dsfont}
\usepackage{xspace}
\usepackage{cancel}
\usepackage{mathtools}
\usepackage{amsthm}
\makeatletter
\@ifundefined{theorem}{\newtheorem{theorem}{Theorem}}{}
\makeatother
\DeclareMathOperator{\Cov}{Cov}
\newcommand{\E}{\mathbb{E}}
\title[Can LLMs Write Mathematics Papers? A Case Study in Reservoir]{Can LLMs Write Mathematics Papers? A Case Study in Reservoir Computing}
\author{Allen G Hart}
\date{Source: arXiv:2509.26550v1 (CC BY 4.0)}
\begin{document}
\maketitle

\noindent\emph{Source and licence.} Can LLMs Write Mathematics Papers? A Case Study in Reservoir Computing. Version arXiv:2509.26550v1 (CC BY 4.0), distributed under the Creative Commons Attribution 4.0 International licence, as recorded in the arXiv licence metadata for this exact version and shown on the abstract page https://arxiv.org/abs/2509.26550v1. Full rights notice: licenses/arxiv-2509.26550-CC-BY-4.0.txt.\par\smallskip

\noindent\textit{Editorial scope.} Counted unit: the Theorem whose source label is \texttt{thm:main} in the article (source0.tex lines 173--190, source label \texttt{thm:main}), delivered together with the article's own complete original proof of it (source0.tex lines 192--239, one \texttt{proof} environment of 48 lines). No mathematics is written, repaired or re-derived: statement and proof are the article's own text line for line. Required context retained uncounted: eq:reservoir\_noise (lines 144--147). Every definition, display and label of those lines is retained unchanged. Omitted from this package and not redistributed: 1--143, 148--172, 191--191, 240--294. Nothing inside a retained excerpt is omitted or altered: every retained body is delivered exactly as the article's lines are written, so no omission receipt of any kind accompanies this package. Citations: the retained excerpts carry no citation command at all, so no bibliography entry of the article is transcribed and no reference is redistributed. Reference adaptation: none was needed and none was made; every reference inside the delivered unit resolves inside it, so no unresolved reference or citation is produced. Printed slips kept literal: no printed word, symbol, hypothesis or number of the counted statement or of its proof was altered. Layout and class: the article's own class is \texttt{article} with packages including amsmath, amssymb, graphicx, hyperref, cite, geometry, amsthm, csquotes, color; the wrapper is the lane's pinned \texttt{amsart} editorial wrapper at 10pt on A4, which restates the theorem-like environments the retained text needs and adds the article's own macro definitions that the retained text needs (\texttt{\textbackslash Cov}, \texttt{\textbackslash E}), with extra wrapper packages (bm, bbm, dsfont, xspace, cancel, mathtools). The delivered numbering is the unit's own. Assets: no retained span contains a figure environment, a graphics-inclusion command, a PDF-page inclusion, a table or a TikZ picture; no image file is copied into this package, the package declares no asset dependency and no image file is redistributed with it. Licence: version arXiv:2509.26550v1 of this article is distributed under the Creative Commons Attribution 4.0 International licence, as recorded in the arXiv licence metadata for this exact version and shown on the abstract page https://arxiv.org/abs/2509.26550v1. A full rights notice is delivered at licenses/arxiv-2509.26550-CC-BY-4.0.txt. Characters: every retained excerpt is pure ASCII, so the delivered engine, XeLaTeX with its default Latin Modern text font, reports no missing character.
\begin{equation}
X_t = AX_{t-1} + C(\omega\phi^t(m) + \xi_t)
\label{eq:reservoir_noise}
\end{equation}

\begin{theorem}[Reservoir Decomposition with Noise]
\label{thm:main}
Consider the reservoir system \eqref{eq:reservoir_noise} with $\rho(A) < 1$. As $t \to \infty$, the reservoir state converges in distribution to:
\begin{equation}
X_\infty = f_{(\phi,\omega,F)}(\phi^t(m)) + Z_\infty
\label{eq:decomposition}
\end{equation}
where:
\begin{enumerate}
    \item $f_{(\phi,\omega,F)}(\phi^t(m))$ is the generalized synchronization function evaluated at the current state
    \item $Z_\infty \sim \mathcal{N}(0, \Sigma)$ is a Gaussian random variable independent of the deterministic dynamics
    \item The covariance matrix $\Sigma$ satisfies the discrete Lyapunov equation:
    \begin{equation}
    \Sigma = A\Sigma A^T + \sigma^2 CC^T
    \label{eq:lyapunov}
    \end{equation}
\end{enumerate}
\end{theorem}

\begin{proof}
We proceed by recursive expansion of the reservoir dynamics. From equation \eqref{eq:reservoir_noise}, we have:
\begin{equation}
X_t = A^t X_0 + \sum_{k=0}^{t-1} A^{t-1-k}C(\omega\phi^{k+1}(m) + \xi_{k+1})
\end{equation}

This can be decomposed into three components:
\begin{equation}
X_t = \underbrace{A^t X_0}_{\text{Initial condition}} + \underbrace{\sum_{k=0}^{t-1} A^{t-1-k}C\omega\phi^{k+1}(m)}_{\text{Deterministic drive}} + \underbrace{\sum_{k=0}^{t-1} A^{t-1-k}C\xi_{k+1}}_{\text{Stochastic component}}
\end{equation}

\textbf{Step 1: Initial condition decay.} Since $\rho(A) < 1$, we have $\|A^t\| \to 0$ exponentially, thus:
\begin{equation}
\lim_{t \to \infty} A^t X_0 = 0
\end{equation}

\textbf{Step 2: Deterministic convergence.} The echo state property ensures that the deterministic sum converges absolutely:
\begin{equation}
\lim_{t \to \infty} \sum_{k=0}^{t-1} A^{t-1-k}C\omega\phi^{k+1}(m) = \sum_{k=0}^{\infty} A^k C\omega\phi^{k+1}(\phi^t(m)) = f_{(\phi,\omega,F)}(\phi^t(m))
\end{equation}
where we used the time-invariance of the infinite sum due to the stationarity of the driving dynamics.

\textbf{Step 3: Stochastic convergence.} Define $Z_t = \sum_{k=0}^{t-1} A^{t-1-k}C\xi_{k+1}$. Since $\xi_k$ are i.i.d. Gaussian with zero mean, $Z_t$ is also Gaussian with:
\begin{equation}
\E[Z_t] = 0
\end{equation}

The covariance matrix is:
\begin{align}
\Cov(Z_t) &= \E[Z_t Z_t^T] = \sum_{k=0}^{t-1} A^{t-1-k}C \E[\xi_{k+1}\xi_{k+1}^T] C^T (A^T)^{t-1-k}\\
&= \sigma^2 \sum_{k=0}^{t-1} A^k CC^T (A^T)^k
\end{align}

Since $\rho(A) < 1$, this sum converges as $t \to \infty$:
\begin{equation}
\lim_{t \to \infty} \Cov(Z_t) = \sigma^2 \sum_{k=0}^{\infty} A^k CC^T (A^T)^k = \Sigma
\end{equation}

where $\Sigma$ satisfies the discrete Lyapunov equation \eqref{eq:lyapunov}, which can be verified by:
\begin{align}
\Sigma &= \sigma^2 \sum_{k=0}^{\infty} A^k CC^T (A^T)^k\\
&= \sigma^2 CC^T + \sigma^2 \sum_{k=1}^{\infty} A^k CC^T (A^T)^k\\
&= \sigma^2 CC^T + A \left(\sigma^2 \sum_{j=0}^{\infty} A^j CC^T (A^T)^j\right) A^T\\
&= \sigma^2 CC^T + A\Sigma A^T
\end{align}

Therefore, $Z_t \xrightarrow{d} Z_\infty \sim \mathcal{N}(0, \Sigma)$ as $t \to \infty$, completing the proof.
\end{proof}

\end{document}
